\textbf{Phần 1. Trắc nghiệm lựa chọn}
\setcounterpageref{ex}{0}

\begin{ex} %[5H5N26]
Hình thang là hình tứ giác có
\choice
{Bốn cạnh bằng nhau}
{\True Một cặp cạnh đối diện song song}
{Hai cặp cạnh đối diện song song}
{Bốn góc vuông}
\loigiai{
Hình thang là hình tứ giác có một cặp cạnh đối diện song song (gọi là hai cạnh đáy).\\
\textbf{Đáp án của câu hỏi là B}
}
\end{ex}

\begin{ex} %[5H5N26]
Hình thang có một cạnh bên vuông góc với hai đáy được gọi là
\choice
{\True Hình thang vuông}
{Hình thang cân}
{Hình chữ nhật}
{Hình bình hành}
\loigiai{
Hình thang có một cạnh bên vuông góc với hai đáy gọi là hình thang vuông.\\
\textbf{Đáp án của câu hỏi là A}
}
\end{ex}

\begin{ex} %[5H5H26]
Diện tích hình thang có độ dài hai đáy lần lượt là $4\text{ cm}$ và $6\text{ cm}$, chiều cao $3\text{ cm}$ là
\choice
{$30\text{ cm}^2$}
{$18\text{ cm}^2$}
{\True $15\text{ cm}^2$}
{$24\text{ cm}^2$}
\loigiai{
Diện tích hình thang là: $\dfrac{(4 + 6) \times 3}{2} = 15\text{ (cm}^2\text{)}$.\\
\textbf{Đáp án của câu hỏi là C}
}
\end{ex}

\begin{ex} %[5H5H26]
Diện tích hình thang có độ dài hai đáy lần lượt là $11\text{ cm}$ và $9\text{ cm}$, chiều cao $8\text{ cm}$ là
\choice
{$160\text{ cm}^2$}
{$40\text{ cm}^2$}
{$99\text{ cm}^2$}
{\True $80\text{ cm}^2$}
\loigiai{
Diện tích hình thang là: $\dfrac{(11 + 9) \times 8}{2} = 80\text{ (cm}^2\text{)}$.\\
\textbf{Đáp án của câu hỏi là D}
}
\end{ex}

\begin{ex} %[5H5H26]
Diện tích hình thang có độ dài hai đáy là $25\text{ cm}$ và $15\text{ cm}$, chiều cao $1\text{ dm}$ là
\choice
{$4\text{ cm}^2$}
{$2\text{ cm}^2$}
{\True $2\text{ dm}^2$}
{$4\text{ dm}^2$}
\loigiai{
Đổi $1\text{ dm} = 10\text{ cm}$. Diện tích hình thang là: $\dfrac{(25 + 15) \times 10}{2} = 200\text{ (cm}^2\text{)} = 2\text{ dm}^2$.\\
\textbf{Đáp án của câu hỏi là C}
}
\end{ex}

\begin{ex} %[5H5H26]
Một hình thang có diện tích $60\text{ cm}^2$, chiều cao $6\text{ cm}$. Tổng độ dài hai đáy của hình thang đó là
\choice
{$10\text{ cm}$}
{$15\text{ cm}$}
{$30\text{ cm}$}
{\True $20\text{ cm}$}
\loigiai{
Tổng độ dài hai đáy là: $60 \times 2 : 6 = 20\text{ (cm)}$.\\
\textbf{Đáp án của câu hỏi là D}
}
\end{ex}

\begin{ex} %[5H5H26]
Một hình thang có diện tích $100\text{ dm}^2$, tổng độ dài hai đáy là $20\text{ dm}$. Chiều cao của hình thang đó là
\choice
{$5\text{ dm}$}
{\True $10\text{ dm}$}
{$15\text{ dm}$}
{$20\text{ dm}$}
\loigiai{
Chiều cao của hình thang là: $100 \times 2 : 20 = 10\text{ (dm)}$.\\
\textbf{Đáp án của câu hỏi là B}
}
\end{ex}

\begin{ex} %[5H5V26]
Một mảnh đất có dạng hình thang vuông với đáy bé $2\text{ m}$, đáy lớn $7\text{ m}$ và chiều cao $13\text{ m}$. Diện tích của mảnh đất đó là
\choice
{$117\text{ m}^2$}
{$58\text{ m}^2$}
{$54\text{ m}^2$}
{\True $58{,}5\text{ m}^2$}
\loigiai{
Diện tích mảnh đất là: $\dfrac{(2 + 7) \times 13}{2} = 58{,}5\text{ (m}^2\text{)}$.\\
\textbf{Đáp án của câu hỏi là D}
}
\end{ex}

\begin{ex} %[5H5V26]
Một mặt bàn đa năng được ghép từ $6$ cái bàn nhỏ hình thang giống nhau có đáy lớn $120\text{ cm}$, đáy bé $60\text{ cm}$ và chiều cao $55\text{ cm}$. Tổng diện tích của mặt bàn đa năng đó là
\choice
{\True $29\,700\text{ cm}^2$}
{$4\,950\text{ cm}^2$}
{$9\,900\text{ cm}^2$}
{$297\text{ m}^2$}
\loigiai{
Diện tích $1$ mặt bàn hình thang là: $\dfrac{(120 + 60) \times 55}{2} = 4\,950\text{ (cm}^2\text{)}$.\\
Diện tích $6$ mặt bàn là: $4\,950 \times 6 = 29\,700\text{ (cm}^2\text{)}$.\\
\textbf{Đáp án của câu hỏi là A}
}
\end{ex}

\begin{ex} %[5H5V26]
Một mảnh đất hình thang có đáy lớn $35\text{ m}$, đáy bé $15\text{ m}$, chiều cao $20\text{ m}$. Người ta mua cỏ để trồng phủ kín mảnh đất với giá $45\,000\text{ đồng/m}^2$. Tổng số tiền mua cỏ là
\choice
{$11\,250\,000\text{ đồng}$}
{\True $22\,500\,000\text{ đồng}$}
{$45\,000\,000\text{ đồng}$}
{$2\,250\,000\text{ đồng}$}
\loigiai{
Diện tích mảnh đất là: $\dfrac{(35 + 15) \times 20}{2} = 500\text{ (m}^2\text{)}$.\\
Số tiền mua cỏ là: $500 \times 45\,000 = 22\,500\,000\text{ (đồng)}$.\\
\textbf{Đáp án của câu hỏi là B}
}
\end{ex}

\begin{ex} %[5H5V26]
Một thửa ruộng hình thang có đáy lớn $120\text{ m}$, đáy bé bằng $\dfrac{2}{3}$ đáy lớn, chiều cao bằng $40\text{ m}$. Diện tích thửa ruộng đó là
\choice
{$4\,800\text{ m}^2$}
{$8\,000\text{ m}^2$}
{\True $4\,000\text{ m}^2$}
{$2\,000\text{ m}^2$}
\loigiai{
Đáy bé thửa ruộng là: $120 \times \dfrac{2}{3} = 80\text{ (m)}$.\\
Diện tích thửa ruộng là: $\dfrac{(120 + 80) \times 40}{2} = 4\,000\text{ (m}^2\text{)}$.\\
\textbf{Đáp án của câu hỏi là C}
}
\end{ex}

\begin{ex} %[5H5V26]
Hình thang có độ dài đáy lớn gấp $3$ lần đáy bé, chiều cao bằng $12\text{ cm}$ và diện tích bằng $96\text{ cm}^2$. Độ dài đáy bé của hình thang đó là
\choice
{\True $4\text{ cm}$}
{$8\text{ cm}$}
{$12\text{ cm}$}
{$16\text{ cm}$}
\loigiai{
Tổng hai đáy là: $96 \times 2 : 12 = 16\text{ (cm)}$.\\
Đáy bé là: $16 : (3 + 1) = 4\text{ (cm)}$.\\
\textbf{Đáp án của câu hỏi là A}
}
\end{ex}

\textbf{Phần 2. Trắc nghiệm đúng sai}
\setcounterpageref{ex}{0}

\begin{ex} %[5H5H26]
Xét tính đúng sai của các khẳng định sau:
\choiceTF
{\True Hai cạnh đáy của hình thang luôn song song với nhau.}
{Hình thang có hai góc vuông thì hai góc vuông đó luôn cùng nằm trên một đáy.}
{\True Diện tích hình thang có hai đáy $14\text{ dm}$ và $6\text{ dm}$, chiều cao $10\text{ dm}$ là $100\text{ dm}^2$.}
{Nếu tăng chiều cao của hình thang lên gấp đôi và giữ nguyên hai đáy thì diện tích tăng gấp bốn lần.}
\loigiai{
\begin{itemize}
    \item Hai cạnh đáy của hình thang luôn song song với nhau.
    \item Hai góc vuông của hình thang vuông nằm ở cùng một cạnh bên (ứng với hai đáy), không phải cùng nằm trên một đáy.
    \item $S = \dfrac{(14 + 6) \times 10}{2} = 100\text{ dm}^2$.
    \item Diện tích tỉ lệ thuận với chiều cao nên tăng chiều cao 2 lần thì diện tích tăng 2 lần (khẳng định tăng 4 lần là sai).
\end{itemize}
\begin{itemize}
    \item a) \textbf{Đúng}.
    \item b) \textbf{Sai}.
    \item c) \textbf{Đúng}.
    \item d) \textbf{Sai}.
\end{itemize}
}
\end{ex}

\begin{ex} %[5H5V26]
Một hình vẽ con thuyền gồm thân thuyền là hình thang có hai đáy $8\text{ cm}$ và $4\text{ cm}$, chiều cao $2\text{ cm}$; hai cánh buồm là hai tam giác vuông có cùng đáy $2\text{ cm}$ và chiều cao $4\text{ cm}$. Xét tính đúng sai của các khẳng định sau:
\choiceTF
{\True Diện tích thân thuyền bằng $12\text{ cm}^2$.}
{\True Diện tích của mỗi cánh buồm là $4\text{ cm}^2$.}
{Tổng diện tích của hai cánh buồm bằng diện tích thân thuyền.}
{\True Tổng diện tích toàn bộ con thuyền là $20\text{ cm}^2$.}
\loigiai{
\begin{itemize}
    \item Diện tích thân thuyền: $\dfrac{(8 + 4) \times 2}{2} = 12\text{ cm}^2$.
    \item Diện tích mỗi cánh buồm: $\dfrac{2 \times 4}{2} = 4\text{ cm}^2$.
    \item Tổng diện tích 2 cánh buồm: $4 + 4 = 8\text{ cm}^2 \neq 12\text{ cm}^2$.
    \item Toàn bộ con thuyền: $12 + 4 + 4 = 20\text{ cm}^2$.
\end{itemize}
\begin{itemize}
    \item a) \textbf{Đúng}.
    \item b) \textbf{Đúng}.
    \item c) \textbf{Sai}: $8\text{ cm}^2 \neq 12\text{ cm}^2$.
    \item d) \textbf{Đúng}.
\end{itemize}
}
\end{ex}

\textbf{Phần 3. Trả lời ngắn}
\setcounterpageref{ex}{0}

\begin{ex} %[5H5H26]
Tính diện tích hình thang có hai đáy lần lượt là $4\text{ cm}$ và $6\text{ cm}$, chiều cao $3\text{ cm}$.
\shortans{15}
\loigiai{
\textbf{Đáp số:} $15$\\
$\dfrac{(4 + 6) \times 3}{2} = 15\text{ (cm}^2\text{)}$.
}
\end{ex}

\begin{ex} %[5H5H26]
Tính diện tích hình thang có hai đáy lần lượt là $6\text{ m}$ và $4\text{ m}$, chiều cao $4\text{ m}$.
\shortans{20}
\loigiai{
\textbf{Đáp số:} $20$\\
$\dfrac{(6 + 4) \times 4}{2} = 20\text{ (m}^2\text{)}$.
}
\end{ex}

\begin{ex} %[5H5H26]
Một hình thang có diện tích $175\text{ cm}^2$, tổng hai đáy là $35\text{ cm}$. Tính chiều cao hình thang đó theo đơn vị xăng-ti-mét.
\shortans{10}
\loigiai{
\textbf{Đáp số:} $10$\\
Chiều cao là: $175 \times 2 : 35 = 10\text{ (cm)}$.
}
\end{ex}

\begin{ex} %[5H5V26]
Mảnh đất hình thang vuông có hai đáy là $2\text{ m}$ và $7\text{ m}$, chiều cao $13\text{ m}$. Tính diện tích mảnh đất đó theo mét vuông. (Nhập kết quả dưới dạng số thập phân).
\shortans{58,5}
\loigiai{
\textbf{Đáp số:} $58{,}5$\\
$\dfrac{(2 + 7) \times 13}{2} = 58{,}5\text{ (m}^2\text{)}$.
}
\end{ex}

\begin{ex} %[5H5V26]
Một mảnh đất hình thang có hai đáy là $35\text{ m}$ và $15\text{ m}$, chiều cao $20\text{ m}$. Tính diện tích mảnh đất đó theo mét vuông.
\shortans{500}
\loigiai{
\textbf{Đáp số:} $500$\\
$\dfrac{(35 + 15) \times 20}{2} = 500\text{ (m}^2\text{)}$.
}
\end{ex}

\begin{ex} %[5H5V26]
Hình thang có diện tích $96\text{ cm}^2$, chiều cao $12\text{ cm}$ và đáy lớn gấp $3$ lần đáy bé. Tính độ dài đáy lớn của hình thang đó theo đơn vị xăng-ti-mét.
\shortans{12}
\loigiai{
\textbf{Đáp số:} $12$\\
Tổng hai đáy: $96 \times 2 : 12 = 16\text{ cm}$. Đáy lớn: $16 : (3 + 1) \times 3 = 12\text{ cm}$.
}
\end{ex}
